\chapter{Materials and Methods}
\label{ch:methods}

\section{Materials}
Membranes, monomers, cleaning chemicals and salts used, with grades and suppliers.

\begin{itemize}
  \item Membranes A and C — end-of-life RO elements (MemBrane Works; one first-pass, one second-pass).
  \item Monomers — \emph{m}-phenylenediamine (MPD), trimesoyl chloride (TMC), \emph{n}-hexane.
  \item Cleaning chemicals — \ce{NaOH}, EDTA, citric acid, \ce{HCl}.
  \item Test salt — \ce{NaCl}.
\end{itemize}


\section{Experimental design}
Two membranes across three cleaning arms give six conditions, with coupons cut into pieces as replicates.

\begin{itemize}
  \item Arm 1 — no clean (IP directly).
  \item Arm 2 — standard clean.
  \item Arm 3 — extensive clean (extended time).
  \item Each coupon cut into 2 pieces for replicates.
  \item Replicated all membranes and arms x2
\end{itemize}

% TODO: insert design schematic / condition table


\section{Cleaning protocols}
Sequential base then acid clean; Arm 3 uses extended duration at matched chemistry.

\begin{itemize}
  \item Base clean — 0.1--0.2\% \ce{NaOH} + 0.8--1\% EDTA.
  \item Rinse / flush.
  \item Acid clean — \ce{HCl} to pH 1--2.
  \item Arm 3 — same chemistry, extended cleaning time.
\end{itemize}


\section{In-situ interfacial polymerisation}
Polyamide is regenerated in place by sequential MPD and TMC dosing through the rig, then heat cured.

\begin{itemize}
  \item MPD (1.5\,wt\%) — 3\,min, waste-first flush.
  \item TMC (0.2\,wt\% in \emph{n}-hexane) — 60\,s.
  \item Heat cure — compressed air via \SI{98}{\celsius} bath ($\approx$\SI{60}{\celsius} at membrane), 5\,min.
  \item Mark coupon orientation.
\end{itemize}


\section{High-pressure crossflow rig and compaction}
Crossflow rig used for compaction and performance testing prior to characterisation.

\begin{itemize}
  \item Compaction — 30\,min at \SI{15}{bar}.
\end{itemize}

% TODO: insert rig diagram


\section{Performance characterisation}
Permeability and salt rejection measured on each piece after compaction.

\subsection{Pure water permeability}
PWP from time-to-volume at known pressure.

\begin{equation}
  A = \frac{V}{A_m \, t \, \Delta P},
  \label{eq:pwp-methods}
\end{equation}

\noindent with $A_m = \SI{0.0139}{\square\metre}$; record $\Delta P$ and temperature.

\subsection{NaCl rejection}
Rejection from feed and permeate conductivities.

\begin{equation}
  R = \left(1 - \frac{\kappa_p}{\kappa_f}\right) \times 100\%,
  \label{eq:rejection-methods}
\end{equation}

\noindent with $\kappa_p$, $\kappa_f$ the permeate and feed conductivities; note feed concentration.


\section{Membrane autopsy}
Surface and chemical analyses on selected pieces to characterise the regenerated layer.

\begin{itemize}
  \item Analyses — FTIR, XPS, SEM-EDS (to confirm).
  \item Coupon-selection criteria — representative of each condition.
\end{itemize}


\section{Safety and waste handling}
Key hazards and disposal routes for the chemicals and procedures used.

\begin{itemize}
  \item Hexane — respirator; never heat with hexane present.
  \item TMC — store under parafilm in glass.
  \item Acid / base — appropriate PPE and handling.
  \item Waste — route per laboratory protocol.
\end{itemize}


\section{Deviations from protocol}
Where execution differed from the Lab Manual; cross-referenced to the Tracker deviation log.

